% Quadratische Gleichungen · p-q-Formel · Prüfungs-Fokus MSA – bank/quadratische-gleichungen/e3.jsonl (zusammenbau v0.6, Prüfungs-Fokus)
\einheitenkopf[e3]{Einheit 3 · Satz vom Nullprodukt}
\zweigzeile{ab Kl. 10 (am Gymnasium ab Kl. 8) · P10 ×7}
\setcounter{aufgabe}{0}

% Anlauf (ohne Prüfkennung)
\begin{aufgabe}{p-q-Formel – Anlauf}
\begin{teile}
\teil $x^2 - x - 12 = 0$ – p und q? p = \leerfeld , q = \leerfeld
\end{teile}
\begin{gleichungsraster}[2]
\gl{x^2 + 6x + 8 = 0} &
\gl{x^2 - 4x - 12 = 0} \\
\end{gleichungsraster}
\end{aufgabe}

\begin{aufgabe}{p-q-Formel – Prüfungsaufgaben}
\begin{teile}
\steil Die Funktion $f$ mit $f(x) = 2x^2 - 2x - 40$ – Nullstellen? \hfill (P20*) \\ x1 = \leerfeld , x2 = \leerfeld
\steil Die Funktion $f$ mit $f(x) = 2x^2 + 16x + 30$ – Nullstellen? \hfill (P20*) \\ x1 = \leerfeld , x2 = \leerfeld
\steil Die Funktion $g$ mit $g(x) = x^2 - 10x + 22$ – Schnittpunkte mit der x-Achse (zwei Stellen)? \hfill (P25*) \\ ( \leerfeld | 0), ( \leerfeld | 0)
\steil Die Funktion $g$ mit $g(x) = x^2 + 4x - 1$ – Schnittpunkte mit der x-Achse (zwei Stellen)? \hfill (P25*) \\ ( \leerfeld | 0), ( \leerfeld | 0)
\steil Parabel $y = -x^2 + 2x + 6$ – an welchen Stellen $x$ ist $y = -2$? \hfill (P23*) \\ x1 = \leerfeld , x2 = \leerfeld
\steil Parabel $y = -x^2 - 6x + 4$ – an welchen Stellen $x$ ist $y = -12$? \hfill (P23*) \\ x1 = \leerfeld , x2 = \leerfeld
\steil Parabel $y = x^2 + 3x - 2$ und Gerade $y = 2x + 4$ – Schnittpunkte? \hfill (P21*) \\ ( \leerfeld | \leerfeld ), ( \leerfeld | \leerfeld )
\steil Parabel $y = x^2 + 4x - 3$ und Gerade $y = x + 1$ – Schnittpunkte? \hfill (P21*) \\ ( \leerfeld | \leerfeld ), ( \leerfeld | \leerfeld )
\steil Parabel $y = (x - 3)^2 - 2$ und Gerade $y = -2x + 19$ – Schnittpunkte? \hfill (P22*) \\ ( \leerfeld | \leerfeld ), ( \leerfeld | \leerfeld )
\steil Parabel $y = (x + 1)^2 - 3$ und Gerade $y = x + 10$ – Schnittpunkte? \hfill (P22*) \\ ( \leerfeld | \leerfeld ), ( \leerfeld | \leerfeld )
\steil Parabel $y = x^2 - 2$ und Gerade $y = 3x + 2$ – Schnittpunkte? \hfill (P24*) \\ ( \leerfeld | \leerfeld ), ( \leerfeld | \leerfeld )
\steil Parabel $y = x^2 - 8$ und Gerade $y = -3x + 2$ – Schnittpunkte? \hfill (P24*) \\ ( \leerfeld | \leerfeld ), ( \leerfeld | \leerfeld )
\steil Zeige, dass $y = (x - 2)^2 - 3$ die Normalform $y = x^2 - 4x + 1$ hat, und berechne die Nullstellen (zwei Stellen). \hfill (P17*) \\ x1 ≈ \leerfeld , x2 ≈ \leerfeld
\steil Zeige, dass $y = (x + 1)^2 - 5$ die Normalform $y = x^2 + 2x - 4$ hat, und berechne die Nullstellen (zwei Stellen). \hfill (P17*) \\ x1 ≈ \leerfeld , x2 ≈ \leerfeld
\end{teile}
\end{aufgabe}

\medskip
\uebersichtskasten{Nullprodukt: Ein Produkt ist null, sobald ein Faktor null ist. Steht rechts null, jeden Faktor einzeln null setzen. \\ x² − 7x = 0 → x·(x − 7) = 0 → x₁ = 0, x₂ = 7 \quad (x + 3)·(x − 9) = 0 → x₁ = −3, x₂ = 9 \\ Ohne Absolutglied immer x ausklammern – nie durch x teilen, sonst geht die Lösung null verloren. \\ Steht rechts nicht null, gilt der Satz nicht: ausmultiplizieren und ordnen. \quad x·(x + 8) = 20 → x² + 8x − 20 = 0 → Formel (Einheit 2)}
